% 第一章：项目介绍
% 负责人：林炜东

\section{项目介绍} \label{sec:intro}

\subsection{项目背景}

随着多核处理器的普及和实时系统复杂度的不断提升，传统的调度算法在多核环境下面临着严重的性能瓶颈。特别是在处理有向无环图(DAG)任务时，缓存的局部性和重用性成为影响系统性能的关键因素。现有的调度算法往往忽略了缓存层次结构对任务执行时间的重要影响，导致系统整体性能无法达到最优。

\subsection{项目目标}

本项目旨在设计并实现一种面向多核实时系统的轻量化缓存感知型调度方法Yat-CASched，主要目标包括：

\begin{enumerate}
    \item 设计基于缓存重用距离的预测模型，准确估算任务在不同核心上的执行时间
    \item 实现高效的在线调度算法，最大化缓存利用率
    \item 开发完整的调度器原型系统，支持DAG任务的实时调度
    \item 通过实验验证算法的有效性和性能优势
\end{enumerate}

\subsection{项目创新点}

\begin{itemize}
    \item 首次提出基于缓存重用距离的DAG任务调度方法
    \item 设计了轻量化的Cache Recency Profile (CRP)预测模型
    \item 实现了高效的在线调度算法，时间复杂度优化至$O(m^2 + m \times (m \times n))$
    \item 提供了完整的系统实现和性能评估框架
\end{itemize}
